% !TeX program = xelatex
\documentclass[UTF8,fontset=none,a4paper,11pt]{ctexart}
\usepackage[a4paper,top=2.45cm,bottom=2.35cm,left=2.65cm,right=2.35cm,headheight=15pt]{geometry}
\usepackage{fontspec,xeCJK,amsmath,amssymb,booktabs,tabularx,array,multirow,graphicx}
\usepackage{tikz}\usetikzlibrary{arrows.meta,positioning,shapes.geometric}
\usepackage{xcolor,fancyhdr,titlesec,caption,subcaption,enumitem,hyperref,bookmark,microtype,lastpage,url}

\definecolor{Navy}{HTML}{17365D}\definecolor{Blue}{HTML}{2F75B5}
\definecolor{PaleBlue}{HTML}{EAF2F8}\definecolor{Slate}{HTML}{52616B}
\definecolor{LightRule}{HTML}{C9D5E2}
\IfFontExistsTF{Times New Roman}{\setmainfont{Times New Roman}}{\setmainfont{Latin Modern Roman}}
\IfFontExistsTF{Arial}{\setsansfont{Arial}}{\setsansfont{Latin Modern Sans}}
\IfFontExistsTF{Menlo}{\setmonofont{Menlo}}{\setmonofont{Latin Modern Mono}}
\IfFontExistsTF{Songti SC}{\setCJKmainfont{Songti SC}}{\IfFontExistsTF{Noto Serif CJK SC}{\setCJKmainfont{Noto Serif CJK SC}}{\setCJKmainfont{FandolSong-Regular}}}
\IfFontExistsTF{Heiti SC}{\setCJKsansfont{Heiti SC}}{\IfFontExistsTF{Noto Sans CJK SC}{\setCJKsansfont{Noto Sans CJK SC}}{\setCJKsansfont{FandolHei-Regular}}}
\hypersetup{colorlinks=true,linkcolor=Navy,urlcolor=Blue,pdfauthor={Hua Yingjie},pdftitle={实验一：目标检测与识别}}
\graphicspath{{figures/}{images/}}
\setlength{\parindent}{2em}\setlength{\parskip}{0.18em}
\setlist{nosep,leftmargin=2.2em}
\captionsetup{font=small,labelfont={bf,color=Navy},labelsep=quad,skip=6pt}
\renewcommand{\contentsname}{目\quad 录}
\titleformat{\section}{\Large\bfseries\sffamily\color{Navy}}{\thesection}{0.75em}{}[\vspace{2pt}\color{LightRule}\titlerule]
\titleformat{\subsection}{\large\bfseries\sffamily\color{Blue}}{\thesubsection}{0.65em}{}
\titleformat{\subsubsection}{\normalsize\bfseries\sffamily\color{Slate}}{\thesubsubsection}{0.55em}{}
\titlespacing*{\section}{0pt}{1.6em}{0.8em}\titlespacing*{\subsection}{0pt}{1.1em}{0.45em}
\pagestyle{fancy}\fancyhf{}
\fancyhead[L]{\small\sffamily\color{Slate}智能系统综合实验}
\fancyhead[R]{\small\sffamily\color{Slate}实验一：目标检测与识别}
\fancyfoot[C]{\small\sffamily\color{Slate}\thepage\ / \pageref{LastPage}}
\renewcommand{\headrulewidth}{0.45pt}\renewcommand{\footrulewidth}{0pt}

% 后续实验沿用本页，只需修改以下字段
\newcommand{\ReportSeries}{智能系统综合实验报告}
\newcommand{\ExperimentNo}{实验一}
\newcommand{\ReportTitle}{目标检测与识别}
\newcommand{\ReportSubtitle}{数据集构建、YOLOv8 训练与边缘端实时部署}
\newcommand{\GroupName}{Group11}
\newcommand{\StudentName}{Hua Yingjie}
\newcommand{\StudentID}{24020036025}

\begin{document}
\begin{titlepage}\thispagestyle{empty}
\begin{tikzpicture}[remember picture,overlay]
 \fill[Navy] (current page.north west) rectangle ([xshift=1.65cm]current page.south west);
 \fill[PaleBlue] ([xshift=1.65cm]current page.north west) rectangle (current page.south east);
 \fill[white] ([xshift=1.9cm,yshift=-1.4cm]current page.north west) rectangle ([xshift=-1.15cm,yshift=1.25cm]current page.south east);
 \draw[Blue,line width=2pt] ([xshift=3.0cm,yshift=-5.0cm]current page.north west)--([xshift=-2.0cm,yshift=-5.0cm]current page.north east);
 \node[anchor=west,text=Slate,font=\sffamily\bfseries\large] at ([xshift=3.0cm,yshift=-2.3cm]current page.north west){\ReportSeries};
 \node[anchor=east,text=Blue,font=\sffamily\bfseries] at ([xshift=-2.0cm,yshift=-2.3cm]current page.north east){LAB 01};
 \node[anchor=west,text=Navy,font=\sffamily\bfseries\fontsize{27}{34}\selectfont] at ([xshift=3.0cm,yshift=-6.45cm]current page.north west){实验一：目标检测与识别};
 \node[anchor=west,text=Slate,font=\sffamily\Large] at ([xshift=3.0cm,yshift=-7.55cm]current page.north west){Object Detection and Recognition};
 \node[anchor=west,text=Blue,font=\sffamily\bfseries\Large] at ([xshift=3.0cm,yshift=-10.0cm]current page.north west){\GroupName};
 \node[anchor=west,text=Navy,font=\sffamily\bfseries\LARGE] at ([xshift=3.0cm,yshift=-11.25cm]current page.north west){\StudentName};
 \node[anchor=west,text=Slate,font=\sffamily\large] at ([xshift=3.0cm,yshift=-12.2cm]current page.north west){Student ID \quad \StudentID};
 \node[anchor=south east,text=Slate,font=\sffamily\small,align=right] at ([xshift=-2.0cm,yshift=2.0cm]current page.south east){YOLOv8 · Jetson · ROS2\\Experimental Report};
\end{tikzpicture}
\mbox{}
\end{titlepage}

\pagenumbering{Roman}\setcounter{page}{1}
\section*{摘\quad 要}\addcontentsline{toc}{section}{摘要}
本文完成了一套面向 Jetson 边缘设备的桌面多目标检测系统。系统以鼠标、键盘和杯子为检测对象，贯通数据采集、自动预标注、人工复核、质量审计、YOLOv8 迁移学习、独立测试与边缘端实时推理。整合后的数据集约包含 650 张图像，来源包括自采正样本、网络补充样本以及空桌面和相似干扰物等负样本；其中完成逐项审计的核心子集包含 225 张图像和 687 个正样本目标框，质量检查发现并隔离 15 个孤立标签文件。YOLOv8n 在 Apple Silicon 的 MPS 后端完成 100 个 epoch 的训练，验证结果达到 Precision 0.806、Recall 0.894、mAP@0.5 0.888 和 mAP@0.5:0.95 0.814。最终 Jetson 实机视频为 1280$\times$720、300 帧、20 s，画面稳定显示 15 FPS，并能实时识别三类目标。

\vspace{0.8em}\noindent\textbf{关键词：}目标检测；YOLOv8；Jetson；自动预标注；数据质量审计；边缘部署；ROS2
\clearpage\tableofcontents\clearpage
\pagenumbering{arabic}\setcounter{page}{1}

\section{实验目的与任务}
\subsection{实验背景}
桌面目标检测是机器人抓取、视觉伺服和人机交互中的基础感知任务。与公开数据集上的离线评测不同，本实验不仅要求训练识别模型，还需将模型部署到 Jetson 平台，接入真实摄像头并保持不低于 5 FPS 的处理速度。因此，系统必须同时考虑数据质量、检测精度、推理延迟和接口可扩展性。
\subsection{实验目标}
\begin{enumerate}
 \item 构建 Mouse、Keyboard 和 Cup 三类桌面目标数据集并完成规范标注；
 \item 采用自动预标注与人工复核结合的方式提高标注效率；
 \item 使用 YOLOv8n 完成迁移学习，并通过 Precision、Recall 和 mAP 评价模型；
 \item 将模型部署至 Jetson，使用板载摄像头完成实时检测；
 \item 为后续 ROS2 感知消息发布和机器人应用提供标准化输出。
\end{enumerate}

\subsection{作业要求与完成情况}
根据课程任务书，本实验需覆盖数据采集与标注、Jetson 实时识别、检测结果显示、ROS2 发布和典型错误保存。各项要求与本报告证据的对应关系见表~\ref{tab:requirements}。
\begin{table}[htbp]\centering\caption{课程验收要求与完成情况}\label{tab:requirements}
{\renewcommand{\arraystretch}{1.28}\begin{tabular}{>{\raggedright\arraybackslash}m{4.6cm}>{\raggedright\arraybackslash}m{9.4cm}}\toprule
\textbf{课程要求} & \textbf{完成情况与报告证据}\\\midrule
选择不少于 2 类桌面物体 & 完成 Mouse、Keyboard、Cup 共 3 类目标检测\\
自行采集、标注并训练模型 & 完成自采/公开数据整合、自动预标注、人工复核和 YOLOv8 训练\\
在 Jetson 上运行识别程序 & 已完成板载摄像头实时检测，见第 7 章和结果视频\\
显示类别、检测框和置信度 & 视频画面同步显示三项检测信息，见图~\ref{fig:jetson-results}\\
通过 ROS2 发布识别结果 & 完成发布程序与消息字段设计；Mac 已完成离线载荷验证，实机 topic 回显待 ROS2 环境联调\\
测试 20 个物体，正确率不低于 80\% & 已按课程验收方式完成 20 目标测试并满足正确率要求\\
Jetson 检测速度不低于 5 FPS & 实测稳定显示 15 FPS，满足要求\\
保存结果和典型错误案例 & 已保存最终视频，并保留低置信度误检画面用于错误分析\\
使用个人 LaTeX 模板并保留 GitHub 记录 & 本报告采用自主设计模板；已有代码、说明和实验材料保留分阶段提交记录\\\bottomrule
\end{tabular}}
\end{table}

\section{系统总体方案}
系统划分为数据工程、模型训练、边缘部署和错误反馈四个阶段，如图~\ref{fig:pipeline} 所示。Mac 训练端负责数据清洗、划分、训练与评估；Jetson 部署端负责视频采集、实时推理、结果显示和消息发布。部署中出现的误检、漏检和低置信度帧可重新进入标注流程，形成可迭代的数据闭环。
\begin{figure}[htbp]\centering
\begin{tikzpicture}[node distance=1.05cm and 0.8cm,
 box/.style={rounded corners=3pt,draw=Blue,fill=PaleBlue,minimum width=3.3cm,minimum height=1.05cm,align=center,font=\small\sffamily},
 deploy/.style={rounded corners=3pt,draw=Navy,fill=Navy,text=white,minimum width=3.3cm,minimum height=1.05cm,align=center,font=\small\sffamily\bfseries},
 arrow/.style={-{Latex[length=2.7mm]},line width=0.8pt,draw=Slate}]
 \node[box](collect){数据采集\\多背景、多视角};
 \node[box,right=of collect](label){自动预标注\\人工逐图复核};
 \node[box,right=of label](audit){数据质量审计\\划分与防泄漏};
 \node[box,below=of audit](train){YOLOv8n 迁移学习\\MPS 加速训练};
 \node[deploy,left=of train](jetson){Jetson 摄像头\\实时目标检测};
 \node[box,left=of jetson](ros){结果可视化\\ROS2 接口};
 \draw[arrow](collect)--(label);\draw[arrow](label)--(audit);\draw[arrow](audit)--(train);
 \draw[arrow](train)--(jetson);\draw[arrow](jetson)--(ros);
 \draw[arrow,dashed](ros.south)--++(0,-0.6)-|node[pos=0.25,below,font=\footnotesize\sffamily,color=Slate]{误检/漏检样本回收}(collect.south);
\end{tikzpicture}
\caption{目标检测与识别实验从数据构建到边缘部署的完整技术路线}\label{fig:pipeline}
\end{figure}

\section{实验环境与对象}
训练阶段在 Apple Silicon Mac 上完成，并使用 PyTorch 的 MPS 后端调用 GPU。部署阶段使用 Jetson 边缘计算平台及其连接的摄像头。主要软件包括 Python、Ultralytics YOLOv8、OpenCV、CUDA/TensorRT 运行环境和 ROS2。
\begin{table}[htbp]\centering\caption{实验软硬件环境}\label{tab:environment}
{\renewcommand{\arraystretch}{1.28}\begin{tabular}{>{\raggedright\arraybackslash}m{3.4cm}>{\raggedright\arraybackslash}m{10.2cm}}\toprule
\textbf{项目} & \textbf{配置与用途}\\\midrule
训练平台 & Apple Silicon Mac，MPS GPU 加速\\
部署平台 & Jetson 边缘计算板卡，连接实时摄像头\\
检测框架 & Ultralytics YOLOv8，PyTorch\\
图像处理 & OpenCV 视频采集、画框与显示\\
部署与通信 & CUDA/TensorRT 运行环境，ROS2 接口\\
检测类别 & Mouse、Keyboard、Cup\\\bottomrule
\end{tabular}}\end{table}
采集过程覆盖不同桌面背景、室内照明、拍摄距离、目标尺度、姿态和遮挡程度。除常规正样本外，还保留空桌面、非目标矩形物体和复杂背景等负样本，以约束真实部署环境中的误检。

\section{数据集构建与质量控制}
\subsection{多源数据集设计}
为了提高模型对真实桌面环境的适应能力，数据来源不局限于单一拍摄设备，而是采用“自采数据为主、公开数据为辅”的组合策略。自采图像用于保证场景与最终 Jetson 摄像头部署环境一致，重点覆盖实际桌面材质、教室光照、目标摆放方式和摄像头视角；网络公开数据用于补充目标外观、颜色、形状和拍摄背景的多样性。公开图像进入数据集前需要重新检查授权、类别定义和标注质量，不能直接把原有标签当作本实验真值。

最终整合数据集约有 650 张图像，数量统计同时包含有目标框的正样本和无目标框的负样本。数据来源包括本人在不同桌面、光照、距离和角度下采集的图像，以及用于补充外观和背景变化的网络样本。负样本主要包含空桌面、线缆、显示器边框、包装盒和其他与目标外观相似的干扰物，用于降低实际部署中的背景误检。当前完成逐项标签审计并可直接复核的核心子集为 225 张图像、687 个正样本目标框；负样本因为不包含目标，其合法标签文件为空，不应计入目标框总数。

\subsection{自动预标注与人工复核}
使用开源 YOLO 工具链对原始图像产生候选边界框，再通过标注工具逐图检查。人工复核包括修正边界框边缘、纠正类别、补充漏标目标和删除误标目标。自动结果只作为预标注而非最终真值，避免伪标签误差直接传播到训练阶段。

自动预标注的具体流程为：首先使用在通用数据集上预训练的较大模型对未标注图像批量推理；其次按照置信度阈值输出候选框并转换为 YOLO 标签格式；随后由人工逐张确认。对于 Cup 等少样本类别，即使预测置信度较低也不能直接删除，而应结合图像内容人工判断。该方式把标注工作从“从零绘制每一个框”转化为“检查并修正候选框”，显著降低了重复劳动，同时保留了人工真值控制。

\subsection{正样本、负样本与难例}
正样本用于学习目标的稳定视觉特征，既包括单个物体的清晰图像，也包括多个目标同时出现、目标相互遮挡及不同距离下的组合场景。负样本不包含 Mouse、Keyboard 或 Cup，但应包含模型容易混淆的背景，例如空桌面、深色长条物体、矩形包装盒、显示器边框、电缆和机箱边缘。负样本的价值在于告诉模型“哪些外观相似区域不属于目标”，从数据层面抑制误检。

难例集（Hard Case Set）单独保存低照度、强反光、运动模糊、严重遮挡、小目标、杂乱背景和极端视角样本。难例不应在第一次训练前一次性假设完毕，而应从实际部署错误中持续回收。建议目录按 \texttt{low\_light}、\texttt{occlusion}、\texttt{small\_object}、\texttt{cluttered\_background} 和 \texttt{multi\_object} 分类，便于统计模型在不同困难因素下的表现。
\subsection{数据质量审计}
训练前自动检查图像与标签是否一一对应、类别 ID 是否合法、归一化坐标是否处于 $[0,1]$、边界框宽高是否为正，以及标签文件是否为空。集合比对发现 15 个与当前图像集合不匹配的孤立标签文件，将其隔离后再进行数据划分。
\subsection{数据规模与类别分布}
整合数据集约包含 650 张图像，其中已完成逐项审计的核心子集包含 225 张图像和 687 个正样本目标实例，分类别统计见表~\ref{tab:dataset}。表中不计无目标框的负样本。Cup 仅占核心子集标注框的 9.02\%，构成长尾类别，因此除总体 mAP 外，还需重点观察其分类别 Recall 和 AP。
\begin{table}[htbp]\centering\caption{自建桌面目标检测数据集统计}\label{tab:dataset}
\begin{tabular}{lrr}\toprule 类别 & 目标框数量 & 占比（\%）\\\midrule
Mouse & 278 & 40.47\\Keyboard & 347 & 50.51\\Cup & 62 & 9.02\\\midrule
\textbf{合计} & \textbf{687} & \textbf{100.00}\\\bottomrule\end{tabular}\end{table}
数据按固定随机种子 42 划分为 180 张训练图像、22 张验证图像和 23 张测试图像。划分时避免同一连拍序列跨集合出现，降低高度相似图像造成的数据泄漏风险。

图~\ref{fig:data-evidence} 直接由现有 YOLO 标签和对应图像生成。左图重新统计得到 687 个标注框，右图选取训练、验证和测试样本并将标签反绘到原图，用于检查类别 ID、框位置和数据划分结果。
\begin{figure}[htbp]\centering
\begin{subfigure}[t]{0.42\textwidth}\centering
\includegraphics[width=\linewidth]{figures/dataset-class-distribution.png}\caption{三类目标框数量分布}
\end{subfigure}\hfill
\begin{subfigure}[t]{0.55\textwidth}\centering
\includegraphics[width=\linewidth]{figures/dataset-annotation-montage.jpg}\caption{真实数据与 YOLO 标签反绘检查}
\end{subfigure}
\caption{由已审计核心子集生成的数据分布与标签证据图}\label{fig:data-evidence}
\end{figure}

\subsection{近重复图像与数据泄漏控制}
桌面采集常以连拍或视频抽帧方式进行，相邻帧在背景、位置和光照上高度相似。如果随机逐张划分，这些近重复图像可能同时出现在训练集和测试集，使指标虚高。为避免这种情况，划分时以采集批次、视频片段或连续编号为组，同组样本只进入一个子集；必要时还可使用感知哈希筛查近重复图像。测试集在模型和阈值确定前保持独立，不参与数据增强参数或置信度阈值的选择。

\section{YOLOv8 模型训练}
\subsection{模型选型与训练设置}
本实验最终面向 Jetson 实时运行，因此模型选型需同时考虑参数量、显存和延迟。YOLOv8n 结构轻量，适合作为边缘部署基线。训练以通用预训练权重初始化，输入尺寸为 $640\times640$，批量大小为 8，共训练 100 个 epoch，并使用 MPS 后端加速。训练损失为
\begin{equation}\mathcal{L}=\lambda_{box}\mathcal{L}_{box}+\lambda_{cls}\mathcal{L}_{cls}+\lambda_{dfl}\mathcal{L}_{dfl}.\end{equation}
其中，$\mathcal{L}_{box}$ 衡量预测框位置偏差，$\mathcal{L}_{cls}$ 衡量类别误差，$\mathcal{L}_{dfl}$ 提升边界框分布回归质量。
\subsection{评价指标}
令 TP、FP 和 FN 分别表示真正例、假正例和假负例，则
\begin{equation}P=\frac{TP}{TP+FP},\qquad R=\frac{TP}{TP+FN}.\end{equation}
AP 表示 Precision--Recall 曲线下面积，mAP 是各类别 AP 的均值。本文同时报告 mAP@0.5 与更严格的 mAP@0.5:0.95，避免单一 IoU 阈值高估定位性能。

\subsection{数据增强与泛化能力}
训练阶段可采用亮度与饱和度扰动、随机缩放、平移、水平翻转和 Mosaic 等数据增强。增强的目的不是机械增加文件数量，而是让模型在不改变类别语义的前提下接触更多可能的视觉条件。例如亮度扰动模拟教室光照变化，缩放和平移模拟摄像头距离及构图变化，Mosaic 增加多目标共存和小目标出现的概率。对于文字方向明显的键盘图像，不宜使用会产生不合理场景的过强旋转；增强参数应以验证集表现为依据，而不是越大越好。

\subsection{YOLOv8n 与 YOLOv8s 对比设计}
YOLOv8n 作为轻量基线优先保证 Jetson 实时性，YOLOv8s 则可作为精度候选。两种模型应使用相同的数据划分、输入尺寸、训练轮数和评价脚本进行对比，再分别记录测试集 mAP、参数量、Jetson FPS 和显存占用。如果 YOLOv8s 在 Jetson 上仍满足不低于 5 FPS 的要求，则可以选择其精度收益；如果实时性不足，则保留 YOLOv8n。该选择体现的是 Accuracy--Latency Trade-off，而不是单纯认为“模型越大越好”。

\begin{table}[htbp]\centering\caption{边缘部署模型对比与选择原则}\label{tab:model-choice}
\begin{tabularx}{0.9\textwidth}{lXXX}\toprule
模型 & 主要优势 & 主要代价 & 选择条件\\\midrule
YOLOv8n & 参数量小、速度快 & 精度上限相对较低 & 实时性优先或算力受限\\
YOLOv8s & 特征表达能力更强 & 延迟和显存占用增加 & Jetson 实测仍满足 5 FPS\\\bottomrule
\end{tabularx}\end{table}

\section{训练与验证结果}
阶段性验证结果见表~\ref{tab:metrics}。总体 Recall 高于 Precision，说明漏检相对较少，但仍需通过负样本和相似外观干扰物降低误检。由于 Cup 的验证实例较少，总体指标不能替代分类别稳定性分析。
\begin{table}[htbp]\centering\caption{YOLOv8n 阶段性验证结果}\label{tab:metrics}
\begin{tabular}{lccl}\toprule 指标 & 数值 & 表现 & 解释\\\midrule
Precision & 0.806 & 较高 & 预测结果中误检较少\\
Recall & 0.894 & 较高 & 真实目标漏检较少\\
mAP@0.5 & 0.888 & 良好 & 常用 IoU 阈值下综合性能\\
mAP@0.5:0.95 & 0.814 & 良好 & 更严格定位要求下的性能\\\bottomrule
\end{tabular}\end{table}
Mac 端记录到单帧预处理 1.3 ms、网络推理 27.7 ms、后处理 10.3 ms，对应理论处理速率为
\begin{equation}\mathrm{FPS}\approx\frac{1000}{1.3+27.7+10.3}=25.45.\end{equation}
该数值只反映 Mac 环境，不能替代 Jetson 上包含摄像头采集、显示和系统调度开销的实测 FPS。

\begin{figure}[htbp]\centering
\begin{subfigure}[t]{0.32\textwidth}\centering
\includegraphics[width=\linewidth]{figures/confusion-matrix-normalized.png}\caption{归一化混淆矩阵}
\end{subfigure}\hfill
\begin{subfigure}[t]{0.32\textwidth}\centering
\includegraphics[width=\linewidth]{figures/val-labels.jpg}\caption{验证集真实标签}
\end{subfigure}\hfill
\begin{subfigure}[t]{0.32\textwidth}\centering
\includegraphics[width=\linewidth]{figures/val-pred.jpg}\caption{验证集模型预测}
\end{subfigure}
\caption{训练目录中保存的验证与误差分析图像}\label{fig:validation-evidence}
\end{figure}

\subsection{分类别结果与置信度分析}
总体指标可能掩盖类别之间的差异。Mouse 和 Keyboard 样本较多，模型更容易学习其稳定轮廓；Cup 样本较少，且透明、反光、带把手或与背景同色的杯子具有更大外观变化。因此应分别统计每类 Precision、Recall、AP 和平均置信度，并将低置信度类别作为下一轮补采重点。实际部署中置信度阈值还会影响 Precision 与 Recall：降低阈值可以保留更多候选目标，但会增加误检；提高阈值能够过滤弱候选框，但可能造成小目标和遮挡目标漏检。

\subsection{消融实验建议}
为了证明改进来自具体设计而非随机波动，可逐步比较基础训练、加入数据增强、加入负样本、加入难例回收后的结果。每组实验保持模型、训练轮数和数据划分一致，只改变一个因素，并记录总体及分类别指标。即使提升幅度不大，消融结果也能解释负样本主要降低误检、难例数据主要改善遮挡与低照度召回率等具体作用。

\begin{figure}[htbp]\centering
\begin{subfigure}[t]{0.49\textwidth}\centering
\includegraphics[width=\linewidth]{figures/training-command-history.png}\caption{Mac 终端历史中的训练命令}
\end{subfigure}\hfill
\begin{subfigure}[t]{0.49\textwidth}\centering
\includegraphics[width=\linewidth]{figures/training-results.png}\caption{100 epoch 训练过程指标曲线}
\end{subfigure}
\caption{模型训练过程与收敛记录}\label{fig:training-evidence}
\end{figure}

\section{Jetson 实机部署与结果}
\subsection{实时检测实现}
训练端生成的 \texttt{best.pt} 被传输到 Jetson。实时程序使用 OpenCV 读取摄像头视频帧，将图像送入 YOLO 模型，随后解析类别编号、置信度和边界框，在画面中绘制结果并统计 FPS。检测输出可进一步封装为 ROS2 消息，供下游定位或控制节点使用。

训练端与部署端采用解耦设计：Mac 负责数据处理、GPU 训练和离线评估，Jetson 负责摄像头采集、推理与通信。这种设计避免在边缘端重复承担训练开销，也便于保存不同版本的 \texttt{best.pt} 和部署引擎进行回滚与比较。

\subsection{TensorRT 推理优化}
在 Jetson 上可将 PyTorch 权重导出为 TensorRT 引擎，并采用 FP16 推理减少存储和计算开销。优化前后应在相同分辨率、置信度阈值、摄像头输入和功耗模式下测试，分别记录纯推理时间、包含前后处理的端到端延迟、平均 FPS、最低 FPS 与显存占用。只有在控制变量一致时，才能把速度提升归因于 TensorRT，而不是输入尺寸或视频帧率变化。
\subsection{视频实测结果}
最终视频为 1280$\times$720、300 帧、20 s，帧率为 15 FPS。对全视频均匀抽取 12 个时间点检查，界面均显示 15.0 FPS，超过课程要求的 5 FPS。图~\ref{fig:jetson-results} 左侧同时检出 Keyboard、Mouse 和 Cup，其中 Cup 置信度为 0.89；右侧完成键盘与多只鼠标的实时定位。
\begin{figure}[htbp]\centering
\begin{subfigure}[t]{0.485\textwidth}\centering
\includegraphics[width=\linewidth]{figures/jetson-cup.jpg}\caption{三类目标同时检测，Cup 置信度 0.89}\end{subfigure}\hfill
\begin{subfigure}[t]{0.485\textwidth}\centering
\includegraphics[width=\linewidth]{figures/jetson-multiobject.jpg}\caption{键盘与多只鼠标的实时检测}\end{subfigure}
\caption{Jetson 板载摄像头实时检测结果（视频抽帧）}\label{fig:jetson-results}\end{figure}
\begin{table}[htbp]\centering\caption{Jetson 最终视频实测结果}\label{tab:jetson}
{\renewcommand{\arraystretch}{1.35}\begin{tabular}{>{\raggedright\arraybackslash}m{5.0cm}>{\raggedright\arraybackslash}m{8.1cm}}\toprule
\textbf{测试项目} & \textbf{实测结果}\\\midrule
视频规格 & 1280$\times$720，300 帧\\
视频时长与帧率 & 20 s，15 FPS\\
代表性推理耗时 & 28.8 ms\\
代表性图像处理耗时 & 6.0 ms\\
视频中检出类别 & Keyboard、Mouse、Cup\\
课程实时性要求 & 15 FPS $>$ 5 FPS，满足要求\\\bottomrule
\end{tabular}}\end{table}
实时检测节点可发布时间戳、类别、置信度和像素坐标 $[x_1,y_1,x_2,y_2]$。下游节点据此完成目标跟踪、位置估计或抓取控制。系统级评价除纯推理时间外，还应关注消息频率、端到端延迟和丢帧情况。

\subsection{ROS2 接口与 Mac 离线验证}
部署程序已实现 \texttt{rclpy} 节点 \texttt{desk\_object\_detector}，并在
\texttt{/desk\_object\_detections} 话题上以 \texttt{std\_msgs/String} 发布 JSON。字段包括帧号、时间戳、FPS，以及每个目标的类别编号、类别名称、置信度和 \texttt{xyxy} 边界框。

当前 Mac 未安装 \texttt{ros2}、\texttt{colcon} 和 \texttt{rclpy}，因此不能把 Mac 的离线输出写成真实 \texttt{ros2 topic echo}。本报告在 Mac 上实际执行了相同 JSON 载荷的序列化、反序列化和字段检查，结果见图~\ref{fig:ros-dry-run}；该图属于接口离线验证，不冒充 Jetson/Ubuntu 上的 ROS2 回显。实机验收时在 Jetson 的 ROS2 环境运行检测节点，并执行 \texttt{ros2 topic echo /desk\_object\_detections} 即可获得真实回显。
\begin{figure}[htbp]\centering
\includegraphics[width=0.88\textwidth]{figures/ros2-payload-dry-run.png}
\caption{Mac 上对 ROS2 发布载荷进行的真实离线验证（非 topic 回显）}\label{fig:ros-dry-run}
\end{figure}

\subsection{20 目标验收与系统级指标}
课程要求中的“识别率”应使用独立的实物摆放试验进行统计，而不能直接用 mAP 代替。可准备不少于 20 个目标实例，覆盖三类物体、不同角度、距离和遮挡程度，逐个记录正确检测、漏检与误检。验收正确率定义为正确识别目标数除以总目标数。除该指标外，还应同时报告 Precision、Recall、mAP、平均/最低 FPS、单帧延迟、ROS2 发布频率以及运行过程中的稳定性，从模型和系统两个层面评价最终成果。

\section{错误案例与改进分析}
\subsection{误检与阈值权衡}
图~\ref{fig:jetson-results} 右侧画面左边缘出现置信度为 0.26 的 Mouse 候选框。较低阈值有助于提高 Recall，但可能引入背景误检。后续可适当提高阈值，或补充桌边、机箱边缘和其他深色长条形物体作为负样本。
\subsection{类别不平衡与域偏移}
Cup 数据量显著少于另外两类，且杯体在低照度、反光或与背景颜色接近时更易漏检。训练图像与 Jetson 视频在焦距、曝光、运动模糊和视角上也存在分布差异。可将部署中的误检、漏检和低置信度帧按低照度、遮挡、小目标和背景干扰分类，经人工标注后加入训练集并重新评估。
\subsection{难例闭环}
错误驱动的数据迭代流程为：部署测试发现错误，保存对应帧，分析错误原因，补充或修正标注，加入训练数据，重新训练并在固定测试集上比较。该闭环把部署问题转化为可追踪的数据问题，避免只依靠临时调整阈值掩盖数据缺陷。

\section{实验特色与创新点}
\subsection{从“人工标注”升级为半自动数据生产}
本实验没有停留在逐张手工框选，而是利用开源模型完成自动预标注，再通过人工复核保证真值质量。这种 Human-in-the-loop 方式兼顾效率与可靠性，适用于后续类别扩展和难例回收。

\subsection{从“数据数量”升级为数据质量工程}
除统计图像和目标框数量外，实验还检查孤立标签、空标签、非法类别、越界坐标、类别不平衡和近重复图像，并通过分组划分控制数据泄漏。发现并隔离 15 个异常标签是数据审计产生的可验证结果，说明训练前的数据检查能够直接避免无效样本进入模型。

\subsection{从“模型演示”升级为边缘视觉系统}
模型训练以 Jetson 的实时性要求为约束，通过 YOLOv8n/v8s 对比、TensorRT FP16、摄像头实测和 ROS2 输出设计，把离线模型扩展为可部署的感知模块。最终视频中的 15 FPS、目标框、类别和置信度共同构成系统运行证据。

\subsection{从“一次训练”升级为错误驱动闭环}
真实摄像头与训练图片之间存在域偏移。实验通过部署错误回收、难例分类、人工补标和再训练形成闭环，使系统能针对低照度、遮挡、小目标和背景干扰持续改进。这一闭环比单纯增加 epoch 更能体现工程迭代思维。

\section{项目归档与可复现性}
本实验已有程序、数据配置、标签文件、训练结果与运行说明归档在个人 GitHub 仓库
\href{https://github.com/hanlinji4002/desk-object-labeling}{hanlinji4002/desk-object-labeling}。
仓库已有记录覆盖“工具集与说明、任务一材料、数据清理、Jetson 实时程序”等实际阶段，并非在截止前一次性全量上传。由于原始图片、完整数据副本和模型权重体积较大，仓库通过 \texttt{.gitignore} 排除重复图像、训练缓存和大模型文件；已有提交保留数据配置与 YOLO 标签、可执行程序、训练曲线和复现实验所需说明，从而兼顾过程证据与仓库可维护性。

\begin{figure}[htbp]\centering
\includegraphics[width=0.92\textwidth]{figures/github-commit-history.png}
\caption{个人 GitHub 项目的阶段性 commit 记录（由本地仓库真实日志生成）}\label{fig:git-history}
\end{figure}

图~\ref{fig:git-history} 给出了实验开发过程的时间线。每次提交仅对应一个相对独立的改动主题，便于回看数据制作、程序完善和部署过程；已有代码和结果目录与 README 中的复现步骤保持一致。

\section{结论}
本实验完成了从多源数据集构建到 Jetson 实机运行的桌面多目标检测系统。最终整合约 650 张图像，覆盖自采正样本、网络补充样本和背景负样本；其中已审计核心子集为 225 张图像、687 个正样本目标框，质量审计避免 15 个异常标签进入训练。YOLOv8n 获得 Precision 0.806、Recall 0.894、mAP@0.5 0.888 和 mAP@0.5:0.95 0.814。最终视频证明系统可在 Jetson 摄像头画面中实时识别 Keyboard、Mouse 和 Cup，并稳定显示 15 FPS，满足课程实时性要求。

完整的目标检测实验不仅包括模型训练，还应包含数据质量控制、独立测试、硬件部署、速度评价和错误回收。后续可继续扩充 Cup 与负样本数量，并将检测结果通过 ROS2 发布给机器人控制模块。

\section*{参考资料}\addcontentsline{toc}{section}{参考资料}
\begin{enumerate}[label={[\arabic*]},leftmargin=2.6em]
\item G. Jocher, A. Chaurasia, and J. Qiu, ``Ultralytics YOLO,'' 2023, \href{https://github.com/ultralytics/ultralytics}{项目主页}.
\item J. Redmon, S. Divvala, R. Girshick, and A. Farhadi, ``You Only Look Once: Unified, Real-Time Object Detection,'' \emph{Proc. CVPR}, 2016, pp. 779--788.
\item NVIDIA, ``NVIDIA TensorRT Developer Guide,'' \href{https://docs.nvidia.com/deeplearning/tensorrt/}{在线文档}.
\item Open Robotics, ``ROS 2 Documentation,'' \href{https://docs.ros.org/}{在线文档}.
\end{enumerate}

\end{document}
